% sections/09-theorem.tex
\section{The theorem}\label{sec:theorem}

This section proves Theorems~\ref{thm:transferintro} and~\ref{thm:openai}. Theorem~\ref{thm:main} and Corollary~\ref{cor:uniform}, with
the rate of Remark~\ref{rem:rate} and the thresholds of Section~\ref{ssec:constants}, are the proof of Theorem~\ref{thm:openai};
Section~\ref{ssec:transferproof} proves Theorem~\ref{thm:transferintro} by the same argument with the data of the certificate left
general.

\subsection{The budget}

Recall $\gs=1-\etas-\frac{61}{4096}-1222\,\eps=4616733319001/10^{14}$, where the last equality is exact arithmetic with
$\etas=93527408868499/10^{14}$ and $\eps=3/10^6$, and $\Cth=72\sum_\nu\norm{\alpha^\nu}^2$ (Appendix~\ref{app:certificate}).

\begin{theorem}[uniform gap on rectangular tori]\label{thm:main}
For every rectangular torus $\Tor$, on the whole Fock space $\Fock$,
\[
\HT^2\ge\KT+\gT\,\HT\ge\gT\,\HT,
\]
where, with $\Ec_m=\Ec_m(\Tor)$, $\lambda_{\min}$ the smallest eigenvalue of the pair Gram matrix and $\dtail=\dtail(\Tor)$,
\begin{equation}\label{eq:gT}
\begin{aligned}
\gT&:=\lambda_{\min}-\etas-\dtail-\Cth\,\Ec_8(1+\Ec_8)-(1+\Ec_{23})^2\bigl(1222\,\eps+2444\,\Ec_{22}\bigr)\\
&\phantom:=\gs-(1-\lambda_{\min})-\Bigl(\dtail-\frac{61}{4096}\Bigr)-\Cth\,\Ec_8(1+\Ec_8)\\
&\phantom{:=\gs}-1222\,\eps\bigl((1+\Ec_{23})^2-1\bigr)-2444\,\Ec_{22}(1+\Ec_{23})^2 .
\end{aligned}
\end{equation}
\end{theorem}

\begin{proof}
By Theorem~\ref{thm:square} and Proposition~\ref{prop:parts},
\[
\HT^2-\KT=\bigl[\LL_2(h^2)-\etas\HT\bigr]+\bigl[\LL_3(X_3(h))-\KT^{(3)}\bigr]+\bigl[\LL_4(W_4(h\otimes h)W_4^*)-\KT^{(4)}\bigr].
\]
By Proposition~\ref{prop:projector}, $h^2\ge\lambda_{\min}h$, and $\LL_2$ preserves order with $\LL_2(h)=\HT$; so the first bracket is at
least $(\lambda_{\min}-\etas)\HT$. The second is bounded by Proposition~\ref{prop:three} and the third by Proposition~\ref{prop:four}.
Finally $\KT\ge0$ and $\HT\ge0$.
\end{proof}

All the constants in \eqref{eq:gT} depend on the torus only through its lattice of periods, and none depends on the
particle number. They are small when the shorter period is large:

\begin{corollary}[uniform bound]\label{cor:uniform}
For $L\ge16$ let
\begin{multline*}
\gl(L):=\gs-\bigl[2\Dlev(L)+\Ups(L,L)\bigr]-\frac32\sum_{\substack{z\ge17\\ z\ \mathrm{odd}}}\frac{(3z-1)(z+1)}{2^z}\,\Xi_z(L)
-\Cth\,\Ec_8(L)(1+\Ec_8(L))\\
-1222\,\eps\bigl((1+\Ec_{23}(L))^2-1\bigr)-2444\,\Ec_{22}(L)(1+\Ec_{23}(L))^2,
\end{multline*}
with $\Ec_m(L)$ and $\Xi_z(L)$ as in Lemma~\ref{lem:Ebounds}. Then $\gT\ge\gl(L)$ for every rectangular torus whose shorter period
is at least $L$, the function $\gl$ is nondecreasing on $[16,\infty)$, and $\gl(L)\to\gs$ as $L\to\infty$.
\end{corollary}

\begin{proof}
By Lemma~\ref{lem:levels}, $1-\lambda_{\min}\le\dG\le\Dlev(L_x)+\Dlev(L_y)+\Ups(L_x,L_y)$. The function $\Dlev$ is decreasing on
$(\sqrt{12},\infty)$, and every term $\abs{\mathfrak f(s_{jk})}=(s_{jk}-1)e^{-s_{jk}/2}$ of $\Ups$ is decreasing in $L_x$ and $L_y$, since off
$\mathcal X$ it has $s_{jk}\ge9/2>3$; so the bound is at most $2\Dlev(L)+\Ups(L,L)$ and is nonincreasing in $L$. The tail term is
bounded by Proposition~\ref{prop:three}, and $\Ec_m(\Tor)\le\Ec_m(L)$ for $m\in\{8,22,23\}$ by Lemma~\ref{lem:Ebounds}(a); all
subtracted terms are nonincreasing in $L$. As $L\to\infty$, $\Dlev(L)\to0$ and $\Ups(L,L)\to0$ (Lemma~\ref{lem:levels}(c)), the tail
excess tends to zero (Proposition~\ref{prop:three}), and $\Ec_m(L)\to0$ (Lemma~\ref{lem:Ebounds}(b)).
\end{proof}

\begin{remark}[the rate of convergence]\label{rem:rate}
As $L\to\infty$,
\[
0\le\gs-\gl(L)=O\bigl(L^4\,2^{-L^2/20}\bigr).
\]
The rate is that of the majorant $\Omega$ of Proposition~\ref{prop:three} for the tail term, and it is not claimed to be sharp.
The same rate holds for the transfer principle, Theorem~\ref{thm:transferintro}: for every planar certificate of the same form,
the bound $\etr(L)$ on $\etr(\Tor)$ obtained in the proof of Theorem~\ref{thm:transfer} is $O(L^4\,2^{-L^2/20})$, with a
constant that depends on the certificate.
\end{remark}

\begin{proof}
$\gs-\gl(L)$ is the sum of the five nonnegative terms subtracted in the definition of $\gl$. By the definition of $\Dlev$ and
Lemma~\ref{lem:levels}(c), $2\Dlev(L)+\Ups(L,L)=O(L^2e^{-L^2/8})$. By Lemma~\ref{lem:Ebounds}(b),
$\Ec_m(L)\le\widehat\Ec_m(L)=O(L^{2m}e^{-L^2/4})$, so the three terms containing $\Ec_8$, $\Ec_{22}$ and $\Ec_{23}$ are
$O(L^{46}e^{-L^2/4})$. By Proposition~\ref{prop:three}, the tail term
$\frac32\sum_{z\ge17,\,z\ \mathrm{odd}}(3z-1)(z+1)2^{-z}\,\Xi_z(L)$ is at most $\sum_{(j,k)\neq0}\Omega(L\sqrt{j^2+k^2})$, which is
$O(L^4\,2^{-L^2/20})$. Since $2^{-L^2/20}=e^{-(\log2)L^2/20}$ and $(\log2)/20<1/8<1/4$, the other four terms are
$o(L^4\,2^{-L^2/20})$. The proof of Theorem~\ref{thm:transfer} bounds the terms of $\etr(\Tor)$ by the same three results, with
$\Ec_{M_c}$, $\Ec_{D_c-1}$ and $\Ec_{D_c}$ in place of $\Ec_8$, $\Ec_{22}$ and $\Ec_{23}$, and the tail by the same majorant (whose
derivation bounds the sum over all $z\ge1$); this gives the second statement.
\end{proof}

\subsection{The constants}\label{ssec:constants}

Table~\ref{tab:constants} lists the terms of $\gl(L)$. The first column after $L$ is the bound $2\Dlev(L)+\Ups(L,L)$ on
$1-\lambda_{\min}$; the binding term is the four-body Bessel term, followed by the tail excess, while the pair-level and
non-orthonormality terms are negligible.

\begin{table}[ht]
\centering\small
\begin{tabular}{@{}rccccc@{}}
\toprule
$L$ & $2\Dlev(L)+\Ups(L,L)$ & $\Ec_8(L)$ & $\Ec_{22}(L)$ & $\Ec_{23}(L)$ & tail excess\\
\midrule
16.5 & $4.5\cdot10^{-13}$ & $2.1\cdot10^{-17}$ & $1.50\cdot10^{-5}$ & $5.73\cdot10^{-5}$ & $5.42\cdot10^{-4}$\\
17   & $5.8\cdot10^{-14}$ & $5.3\cdot10^{-19}$ & $1.13\cdot10^{-6}$ & $4.73\cdot10^{-6}$ & $1.93\cdot10^{-4}$\\
17.5 & $7.2\cdot10^{-15}$ & $1.2\cdot10^{-20}$ & $6.95\cdot10^{-8}$ & $3.19\cdot10^{-7}$ & $8.38\cdot10^{-5}$\\
18   & $8.2\cdot10^{-16}$ & $2.2\cdot10^{-22}$ & $3.56\cdot10^{-9}$ & $1.77\cdot10^{-8}$ & $4.22\cdot10^{-5}$\\
19   & $9.0\cdot10^{-18}$ & $5.2\cdot10^{-26}$ & $5.45\cdot10^{-12}$ & $3.16\cdot10^{-11}$ & $4.05\cdot10^{-6}$\\
20   & $7.6\cdot10^{-20}$ & $7.1\cdot10^{-30}$ & $4.18\cdot10^{-15}$ & $2.79\cdot10^{-14}$ & $4.95\cdot10^{-7}$\\
\bottomrule
\end{tabular}

\medskip
\begin{tabular}{@{}rcccc@{}}
\toprule
$L$ & $\Cth\Ec_8(1+\Ec_8)$ & $1222\,\eps\bigl((1+\Ec_{23})^2-1\bigr)$ & $2444\,\Ec_{22}(1+\Ec_{23})^2$ & $\gl(L)$\\
\midrule
16.5 & $6.0\cdot10^{-15}$ & $4.2\cdot10^{-7}$ & $3.67\cdot10^{-2}$ & $0.0089238941$\\
17   & $1.5\cdot10^{-16}$ & $3.5\cdot10^{-8}$ & $2.75\cdot10^{-3}$ & $0.0432229058$\\
17.5 & $3.3\cdot10^{-18}$ & $2.3\cdot10^{-9}$ & $1.70\cdot10^{-4}$ & $0.0459136282$\\
18   & $6.2\cdot10^{-20}$ & $1.3\cdot10^{-10}$ & $8.70\cdot10^{-6}$ & $0.0461164613$\\
19   & $1.5\cdot10^{-23}$ & $2.3\cdot10^{-13}$ & $1.33\cdot10^{-8}$ & $0.0461632655$\\
20   & $2.0\cdot10^{-27}$ & $2.0\cdot10^{-16}$ & $1.02\cdot10^{-11}$ & $0.0461668382$\\
\bottomrule
\end{tabular}
\caption{The terms of the uniform bound $\gl(L)$ of Corollary~\ref{cor:uniform}. The tail excess is
$\frac32\sum_{z\ge17\ \mathrm{odd}}(3z-1)(z+1)2^{-z}\,\Xi_z(L)$, and $\gs=0.04616733319001$.}
\label{tab:constants}
\end{table}

\paragraph{Thresholds.} The computed function $\gl$ crosses $0$ at $L=16.4562$, $1/25$ at $L=16.8566$, and $\gs-10^{-5}$ at
$L=18.6843$ (magnetic lengths, to the precision shown). Hence $\gl(L)>0$ for $L\ge16.457$, $\gl(L)>1/25$ for $L\ge16.857$, and
$\gl(L)\ge\gs-10^{-5}$ for $L\ge18.685$; by Corollary~\ref{cor:uniform} the same holds for $\gT$ on every rectangular torus whose
shorter period is at least these values.

\begin{remark}[on the evaluation of the constants]\label{rem:evaluation}
Every quantity in $\gl(L)$ is given by an explicit expression: $\Dlev$ is elementary, $\Ups(L,L)$ and $\Xi_z(L)$ are lattice sums of
elementary functions with Gaussian decay, the sum over odd $z$ converges geometrically (Proposition~\ref{prop:three}), and
$\mu_m(r)$ is the spectral norm of an explicit $(m+1)\times(m+1)$ matrix whose entries are Laguerre polynomials times Gaussians
(Lemma~\ref{lem:bargmann}(iii)). The reduction from the supremum over all rectangular tori to the evaluation at the radii $L\sqrt{j^2+k^2}$ is
rigorous by Lemma~\ref{lem:Ebounds}(a), whose proof shows that for $r\ge16$ every entry of $\Dis^{(m)}(r/\sqrt2)$, $m\le23$, has a fixed
sign and a modulus decreasing in $r$, so that the spectral norm is a monotone function of the entries. What remains is a
numerical evaluation. The entries of Table~\ref{tab:constants} were computed in $60$-digit floating-point arithmetic from these
expressions, with $\gs$, $61/4096$ and $\etas$ as exact rationals, the displacement entries from the Laguerre representation, and
the spectral norms in double precision cross-checked against $60$-digit eigenvalue computations (relative agreement
$3\cdot10^{-16}$); the lattice sums and the series over $z$ were truncated where the omitted terms are far below the displayed
digits. A second, independent implementation with different truncations reproduces every entry to the digits it reports (four
to seven significant figures for $\gl$). The values are therefore
floating-point evaluations of rigorous closed-form bounds, not interval-arithmetic certificates, and the thresholds were located
by bisection of the computed function.
\end{remark}

\subsection{Consequences}

\begin{corollary}[spectrum and kernel]\label{cor:spectrum}
Let $\Tor$ be a rectangular torus with $\gT>0$; by Corollary~\ref{cor:uniform} this holds whenever the shorter period is at least
$16.457$.
\begin{enumerate}[(a)]
\item For every $N$, $\spec\bigl(\HT|_{\bigwedge^N\h}\bigr)\subset\{0\}\cup[\gT,\infty)$.
\item At $\Nphi=3N$, $\HT|_{\bigwedge^N\h}\ge\gT(\one-P_N)$, where $P_N$ is the orthogonal projection onto the three-dimensional span of the
Laughlin states $\Psi_0,\Psi_1,\Psi_2$ of Theorem~\ref{thm:kernel}, and $\HT|_{\bigwedge^{N'}\h}\ge\gT$ for every $N'>N$.
\item If $\dG<1$, the projector form satisfies $H'^2\ge(\gT/\lambda_{\max})H'$.
\end{enumerate}
In particular, on every rectangular torus whose shorter period is at least $16.857$, the gap above the three-dimensional kernel at
$\Nphi=3N$ is larger than $1/25$, uniformly in $N$ and in the aspect ratio.
\end{corollary}

\begin{proof}
(a) $\HT\ge0$ preserves the particle number, and an eigenvalue $\lambda$ satisfies $\lambda^2\ge\gT\lambda$ by Theorem~\ref{thm:main}.
(b) By (a) and the spectral theorem, $\HT\ge\gT(\one-P_{\ker})$ on each sector; Theorem~\ref{thm:kernel} identifies the kernel at
$N$ particles and Proposition~\ref{prop:overfill} shows that it is zero for $N'>N$. (c) By Proposition~\ref{prop:projector},
$H'\ge\lambda_{\max}^{-1}\HT\ge(\gT/\lambda_{\max})(\one-P_{\ker})$ with $\ker H'=\ker\HT$.
\end{proof}

At the square aspect ratio the shorter period is $L=\sqrt{2\pi\Nphi}$, so $\gl>0$ for $\Nphi\ge44$ and $\gl>1/25$ for $\Nphi\ge46$.
At filling one third, $\Nphi=3N$: for $N=15$ ($L=16.815$) the bound is $\gl\ge0.0385$, and for $N=16$ ($L=17.366$) it is
$\gl\ge0.0456>1/25$. On elongated tori the shorter period is smaller at the same $N$, and correspondingly larger systems are
needed.

\subsection{The transfer principle}\label{ssec:transferproof}

It remains to prove Theorem~\ref{thm:transferintro}, for a planar certificate of the same form, with the inequalities
\eqref{eq:certgeneral} and the constant $\gamma$ of the definition in Section~\ref{ssec:transfer}. For the certificate
of~\cite{OAI:A-Fock-space-inequality-and-the-Laughlin-spectral-gap-September-24-2026}, $P_c=7$, $M_c=8$, $Z_c=15$, $D_c=23$,
$\delta_\infty(15)=61/4096$, $N_4=1222$ and $\gamma=\gs$. We restate the theorem in its full form.

\begin{theorem}[the transfer principle; Theorem~\ref{thm:transferintro} restated]\label{thm:transfer}
Given a planar certificate of the same form with constant $\gamma$, for every rectangular torus, on $\Fock$,
\[
\HT^2\ge(\gamma-\etr(\Tor))\,\HT,
\]
where, with $\Ec_m=\Ec_m(\Tor)$ and $\sigma_z=\sigma_z(\Tor)$,
\begin{multline*}
\etr(\Tor):=(1-\lambda_{\min})+\frac32\sum_{\substack{z>Z_c\\ z\ \mathrm{odd}}}\frac{(3z-1)(z+1)}{2^z}\sigma_z
+(M_c+1)(P_c+1)\sum_\nu\norm{\alpha^\nu}^2\,\Ec_{M_c}(1+\Ec_{M_c})\\
+\eps N_4\bigl((1+\Ec_{D_c})^2-1\bigr)+2N_4\,\Ec_{D_c-1}(1+\Ec_{D_c})^2 .
\end{multline*}
For every rectangular torus whose shorter period is at least $L$, $\etr(\Tor)\le\etr(L)$ with $\etr(L)\to0$ as $L\to\infty$.
In particular, if $\gamma>0$, then for every $\gamma'<\gamma$ there is $L_0$ such that $\HT^2\ge\gamma'\HT$ on the whole Fock space of every
rectangular torus whose shorter period is at least $L_0$, uniformly in the particle number and the aspect ratio.
\end{theorem}

\begin{proof}
The proofs of Proposition~\ref{prop:parts}, Proposition~\ref{prop:three}, Lemma~\ref{lem:bessel}, Proposition~\ref{prop:four} and
Theorem~\ref{thm:main} use the specific certificate only through the data of the definition in Section~\ref{ssec:transfer}, and go
through verbatim with the following changes. In Proposition~\ref{prop:parts}, the modes in $(i,k)$ and $j$ are at most $M_c$, each pair label $p\le P_c$
occurs with at most $M_c+1$ values of $j$, and $\frac1\pi\int\sum_{p\le P_c}\norm{B(\ut_p)\Psi}^2=(P_c+1)\ip\Psi{\HT\Psi}$, which gives the
constant $(M_c+1)(P_c+1)\sum_\nu\norm{\alpha^\nu}^2$ with $\Ec_{M_c}$; the three-body blocks are $z\le Z_c$, the four-body blocks
$D\le D_c$, with mode indices at most $D_c$ (whence $\Ec_{D_c}$) and $\sum_{D\le D_c}n_D=N_4$. In Proposition~\ref{prop:three} the near
part uses the first family in \eqref{eq:certgeneral} (non-strict inequalities suffice, and the blocks $z\in\{0,1,3\}$ vanish), the
even blocks $z>Z_c$ are positive and the odd ones are bounded as before. In Lemma~\ref{lem:bessel} the pair-centre channels are
$d\le D_c-1$, whence $\Ec_{D_c-1}$, and Proposition~\ref{prop:four} gives $-(1+\Ec_{D_c})^2(\eps N_4+2N_4\Ec_{D_c-1})\HT$. Collecting the
terms as in \eqref{eq:gT} gives the first claim. For the second, $1-\lambda_{\min}\le2\Dlev(L)+\Ups(L,L)$ (proof of
Corollary~\ref{cor:uniform}), the tail is bounded by Proposition~\ref{prop:three}, and $\Ec_m(\Tor)\le\widehat\Ec_m(L)$ for
$L^2\ge2m$ by Lemma~\ref{lem:Ebounds}(b); each bound tends to zero.
\end{proof}
